\documentclass[11pt,twoside]{article}

\usepackage[
  a4paper,
  margin=1.5cm,
  includehead, includefoot
]{geometry}

\usepackage[utf8]{inputenc}
\usepackage[brazil,portuges]{babel}
\usepackage{csquotes}

\usepackage[
  backend=biber,
  style=abnt,
  language=auto,
]{biblatex}

  \addbibresource{../../refs/referencias.bib}

\usepackage{soul}
\usepackage{xcolor}
\usepackage{framed}

\usepackage{tikz}

  \usetikzlibrary{shapes.geometric, positioning, decorations.pathreplacing}

\usepackage{kpfonts}
\usepackage[T1]{fontenc}

\usepackage{tabularx}

  \newcolumntype{C}{>{\centering\arraybackslash}X}

\usepackage{graphicx}
\usepackage{multirow}
\usepackage{multicol}
\usepackage{hhline}
\usepackage{booktabs}
\usepackage[inline]{enumitem}
\usepackage{xfrac}
\usepackage{mathtools}

  \newcommand{\Binary}{\ensuremath{\mathbb{B}}}
  \newcommand{\Natural}{\ensuremath{\mathbb{N}}}
  \newcommand{\Integer}{\ensuremath{\mathbb{Z}}}
  \newcommand{\Rational}{\ensuremath{\mathbb{Q}}}
  \newcommand{\Real}{\ensuremath{\mathbb{R}}}

  \DeclarePairedDelimiter{\floor}{\lfloor}{\rfloor}
  \DeclarePairedDelimiter{\ceil}{\lceil}{\rceil}
  \DeclarePairedDelimiter{\abs}{\lvert}{\rvert}

\usepackage{amsthm}

  \newtheorem{theorem}{Teorema}[section]
  \newtheorem{proposition}[theorem]{Proposição}
  \newtheorem{lemma}[theorem]{Lema}
  \newtheorem{corollary}[theorem]{Corolário}
  \theoremstyle{definition}
  \newtheorem{definition}{Definição}[section]

\usepackage{multicol}

  \newcommand{\finishcols}{\vfill\vphantom{.}}

\usepackage[portuguese,algoruled,vlined,noend]{algorithm2e}

  \SetKwInput{Algoritmo}{Algoritmo}
  \SetKw{Retorne}{retorne}
  \SetKwFor{Para}{para}{faça}{fimpara}
  \SetKw{Escreva}{escreva}
  \DontPrintSemicolon

\usepackage{hyperref}
\usepackage[brazil]{varioref}
\usepackage[brazilian]{cleveref}

  % !FIXME: Corrigir nome automático dos ambientes a seguir
  \crefname{definition}{Definição}{Definições}
  \crefname{lemma}{Lema}{Lemas}
  \crefname{corollary}{Corolário}{Corolários}

  % !FIXME: Mover bloco a seguir para final do preâmbulo
  \def\disciplinacodigotext{QXD0115}
  \def\disciplinanometext{Estrutura de Dados Avançada}
  \def\semestretext{2022.1}
  \def\provatype{Notas de Aula}
  \def\provatitle{Árvores Rubro-Negras são balanceadas}
  \def\provatext{\provatype{} \ifx\provatype\empty {} \else {--} \fi \provatitle{}}
  \def\provacodtext{}
  \def\professorabbrv{Prof.}
  \def\professornometext{Arthur Araruna}

\usepackage{fancyhdr}

  \pagestyle{fancy}
  \fancyhead{}
  \lhead{\disciplinacodigotext}
  \chead{\provatext}
  \rhead{\semestretext}
  \fancyfoot{}
  \fancyfoot[LE,RO]{\colorbox{black!20}{\thepage}}

% * NOTE: Redefinição feita para evitar confusão com exponenciações.
\renewcommand{\thefootnote}{(\roman{footnote})}

\begin{document}
\noindent
\begin{minipage}{\textwidth}
	\begin{tabularx}{\textwidth}{cXc}
		\multirow{4}{*}{\includegraphics[height=1.8cm]{../brasaoufc.pdf}} & {\bf UNIVERSIDADE FEDERAL DO CEARÁ}         & \multirow{3}{*}{} \\
		                                                                  & {\bf Campus de Quixadá}                     &                   \\
		                                                                  & \professorabbrv\enspace\professornometext   &                   \\
		                                                                  & \disciplinacodigotext - \disciplinanometext & {}                \\
	\end{tabularx}
\end{minipage}

\part*{\provatitle}

\section{Conceitos}

Vamos relembrar o nossos conceitos que serão necessários para o resultado a ser demonstrado:

\begin{definition}
	O \emph{tamanho} de uma árvore é a sua quantidade de nós.
\end{definition}

\begin{definition}
	Uma árvore binária \( T \), com \( n \) nós, tal que \( h(T) = O(\lg n) \) é dita \emph{balanceada}.
\end{definition}

\begin{definition}
	As sub-árvores vazias em uma árvore binária são consideradas nós virtuais (existem apenas conceitualmente) chamados de \emph{nós externos}, sendo os demais chamados de \emph{nós internos} (os que realmente guardam informações).

	Nessa ótica, toda árvore binária é estrita e as folhas são exatamente os nós externos.
\end{definition}

\begin{definition}
	Uma árvore binária estrita \( T \) é \emph{Rubro-negra} quando é possível colorir os nós de \( T \) com as cores \emph{vermelho} ou \emph{preto} de forma que todas as propriedades a seguir são válidas:
	\begin{enumerate}[label=(\alph*)]
		\item todo nó é ou vermelho ou preto.
		\item a raiz é preta.
		\item todos os nós externos são pretos.
		\item se um nó é vermelho, ambos os seus filhos são pretos.
		\item para todo nó \( n \), todos os caminhos descendentes de \( n \) às folhas na sua sub-árvore contém exatamente a mesma quantidade de nós pretos.
	\end{enumerate}
\end{definition}

\begin{definition}
	A \emph{altura em negros} (ou \emph{altura-negra}) de um nó \( n \) de uma árvore Rubro-Negra, denotada por \( hn(n) \), é a quantidade de nós pretos em um caminho descendente de \( n \) a uma das suas folhas.

	A altura em negros de uma árvore \( T \), denotada por \( hn(T) \), é a altura em negros de sua raiz.
\end{definition}

\section{Demonstração}

\begin{framed}
	{\bfseries Observação:}

	A argumentação abaixo não é baseada na que consta em~\textcite{Szwarcfiter2010} nem em~\textcite{Cormen2009}.
\end{framed}

\begin{lemma} \label{res:hn-min-lgn}
	Uma árvore Rubro-Negra \( T \) de tamanho \( n \) \emph{mínimo} é tal que \( hn(T) = O(\lg n) \).
\end{lemma}
\begin{proof}
	Seja \( T \) uma Árvore Rubro-Negra de altura-negra \( hn(T) = k \) para um \( k \) arbitrário e considere que \( T \) tenha o menor tamanho \( n \) dentre todas as Rubro-Negras com essa altura-negra.

	Observe que em \( T \) não podem existir nós vermelhos, pois como esses nós não são contados no valor de \( k \), poderíamos substitui-los junto com seus filhos por um nó preto e reduzir o tamanho da árvore em duas unidades para cada nó vermelho sem alterar a altura-negra. % TODO: fazer figura ilustrando.

	Dessa forma, pela definição de Árvores Rubro-Negras e de altura-negra, como todos os nós devem ser pretos temos que essa árvore deve ser perfeita e, certamente, \( k = hn(T) = O(\lg n) \)\footnote{Como \( k \) é arbitrário, temos que esse resultado é válido para todas as árvores.}.
\end{proof}

\begin{corollary}
	Uma árvore Rubro-Negra \( T \) qualquer de tamanho \( n \) \emph{qualquer} é tal que \( hn(T) = O(\lg n) \).
\end{corollary}
\begin{proof}
	Considere \( T \) e \( T^\prime \) como árvores Rubro-Negras de mesma altura-negra, sendo \( T \) de tamanho \( n \) mínimo e \( T^\prime \) de tamanho \( n^\prime \). Certamente \( n \leq n^\prime \) e por isso \( \lg n \leq \lg n^\prime \implies \lg n = O(\lg n^\prime) \). Como, pelo Lema~\ref{res:hn-min-lgn}, \( hn(T) = O(\lg n) \), temos que \( hn(T^\prime) = hn(T) = O(\lg n) = O(\lg n^\prime) \).
\end{proof}

\begin{lemma}
	Para uma árvore Rubro-Negra \( T \) qualquer, \( h(T) \leq 2 hn(T) \).
\end{lemma}
\begin{proof}
	Considere \( T^\prime \) Rubro-Negra com altura-negra \( k \), com \( k \) arbitrário, e altura máxima \( h_{max} \). Perceba que uma Rubro-Negra não pode ter altura máxima quando o seu caminho mais comprido da raiz a alguma folha\footnote{Esse é o caminho que determina a altura da árvore.} possui ao menos um nó preto que seja pai de outro nó preto, pois podemos adicionar um nó vermelho entre esses\footnote{junto com alguns outros nós para completar a outra sub-árvore desse nó vermelho.} e aumentar a altura da árvore sem alterar sua altura-negra.

	Assim, em \( T^\prime \) o maior caminho certamente é constituído de uma alternância entre nós pretos e vermelhos. A quantidade de nós nesse caminho é igual a \( h_{max} \) porém apenas metade (aproximadamente) deles são contados para \( k \). Assim, temos que \( h_{max} = 2 k \). Então, dada uma Rubro-Negra \( T \) qualquer com \( hn(T) = k \), temos que \( h(T) \leq h(T^\prime) = h_{max} \), pela escolha de \( T^\prime \) e daí concluímos que \( h(T) \leq 2k = 2hn(T) \)\footnote{Como \( k \) é arbitrário, temos que esse resultado é válido para todas as árvores.}.
\end{proof}

\begin{theorem}
	Árvores Rubro-Negras são balanceadas.
\end{theorem}
\begin{proof}
	Como para qualquer Rubro-Negra \( T \) temos \( h(T) \leq 2 hn(T) \) e \( hn(T) = O(\lg n) \), temos que \( h(T) = O(\lg n) \).
\end{proof}

\printbibliography

\end{document}